\section{Diagram Families}
\label{sec:families}

A \texttt{DiagramKind} is one registered module. They fall into a Mermaid-compatible
group and four Triton-native families. The detector recognises the Mermaid headers
and the Triton-native ones alike (\texttt{src/frontend/detect.ts}).

\subsection{Mermaid-compatible}

These accept Mermaid syntax and render it unchanged, so existing diagrams migrate
with zero edits:

\begin{lstlisting}[language=triton]
flowchart LR
  start[Start] --> validate{Valid?}
  validate -->|yes| process[Process]
  validate -->|no|  reject[Reject]
\end{lstlisting}

\ourfig[0.6\linewidth]{flowchart}{A Mermaid flowchart, rendered unchanged through
the \texttt{flowchart} module.}

The set spans node-link (flowchart), UML/software (sequence, class, state, ER),
project (gantt, timeline), and chart kinds. The charts (pie, xychart, quadrant,
radar) are \emph{four separate modules}, not a single ``grammar of graphics'' ---
each owns the IR that fits it.

\subsection{Tree family}

A single decorated \texttt{tree} engine underlies the family. Its IR is a
\emph{flat, id-referenced node list} --- not recursively nested nodes --- where
each node carries a small decoration vocabulary (shape, colour class, corner
badge, sub-label, edge label). The tidy layout from \S\ref{sec:kernels} places it.

On top of that engine sit \emph{value-driven} front-ends. Their source is a list
of operations, and the builder computes a valid structure --- you cannot author an
invalid one:

\begin{lstlisting}[language=triton]
avl insert 50 30 70 20 40 60 80 10
\end{lstlisting}

\ourfig[0.72\linewidth]{avl}{An AVL tree built from the \texttt{insert} list above.
The balance-factor badge on each node is computed by running real insertion with
rotations --- the structure cannot be authored unbalanced.}

\noindent renders a balanced AVL tree with a balance-factor badge on every node,
because the builder runs real insertion with rotations. The siblings work the
same way: \texttt{rbtree} (red/black colouring), \texttt{btree order N} (multi-key
strip nodes), \texttt{radix} (prefix-compressed tries with substring edge labels),
\texttt{segtree over [\dots] reduce sum} (range-annotated nodes), \texttt{heap max}
(sift-up). \texttt{plan} reuses the same engine for query-plan trees, inferring
operator colour from the node label.

\ourfig[0.78\linewidth]{radix}{A radix (Patricia) trie over a set of words: shared
prefixes are compressed onto single labelled edges.}

\subsection{Struct / memory family}

These exercise the cell-strip and slot-aware pointer kernels:

\begin{lstlisting}[language=triton]
array
  title int[] nums
  cells 5 8 13 21 34
  index
  ptr i -> 2
\end{lstlisting}

\ourfig[0.7\linewidth]{array}{An array with an index row and a named pointer into
a cell.}

\texttt{linkedlist} draws a \texttt{head} chain of \texttt{[value|next]} nodes
terminating in a null slash; \texttt{memory} lays out named regions (e.g.\ a stack
and a heap) and routes a \emph{cross-region pointer} from a variable slot to a heap
object via \texttt{connectSlots}; \texttt{page} draws a slotted storage page whose
line pointers indirect to tuples.

\ourfig[0.8\linewidth]{memory}{A memory model: a variable slot in one region holds a
pointer that crosses into a heap object in another, routed by the slot-aware
connector.}

\subsection{Topology}

A cost-weighted node graph. Edges are coloured and dashed by the tier their weight
falls into, a legend is generated from the declared scale, and nodes may be nested
inside group containers:

\begin{lstlisting}[language=triton]
topology
  title NUMA interconnect
  costs ns
    tier local 90 #27ae60
    tier hop1  140 #2f80ed
    tier hop2  200 #e2574c 5 4
  node N0 : Node 0
  node N1 : Node 1
  N0 -- N1 : 140
\end{lstlisting}

\ourfig{numa}{A NUMA topology with nested node groups: links are coloured and
dashed by the cost tier their weight falls into, and the legend is generated from
the declared scale.}
